% Unidad U006F. Inventario operatorio completo del Topological Phase Kernel.

\chapter{Inventaire opératoriel complet et chaîne causale du Topological Phase Kernel}
\label{chap:inventario-operatorio-completo-tpk}

Le Topological Phase Kernel n’est ni une réduction modulo trois ni une horloge à
cent huit positions. C’est une catégorie d’états enrichis équipée
d’opérateurs de sélection, de transport, de lecture, de mémoire, de périodicité,
de prolongement et de publication. Pour empêcher une abréviation de remplacer
l’objet, chaque opérateur est enregistré selon son type et l’information qu’il
conserve.

\section{Fiche typée d'un opérateur}

\begin{definition}[Fiche opératoire]
\label{def:u006f-ficha-operatoria}
Une réalisation canonique du TPK est un sextuplet
\begin{equation}
 \mathcal O=
 (D_{\mathcal O},C_{\mathcal O},f_{\mathcal O},
 I_{\mathcal O},L_{\mathcal O},M_{\mathcal O}),
 \label{eq:u006f-ficha}
\end{equation}
où :
\begin{enumerate}
 \item \(f_{\mathcal O}:D_{\mathcal O}\to C_{\mathcal O}\) est sa règle ;
 \item \(I_{\mathcal O}\) est la famille des relations conservées ;
 \item \(L_{\mathcal O}\) déclare l’information publiée ou perdue ;
 \item \(M_{\mathcal O}\) énumère la mémoire qui demeure dans l’état
 enrichi.
\end{enumerate}
\end{definition}

Deux réalisations ne représentent le même opérateur que s’il existe une
équivalence typée qui entrelace les règles et les invariants. Partager une période,
une cardinalité ou une sortie visible ne suffit pas.

\section{Huit classes structurelles}

L’inventaire est partitionné en :
\begin{enumerate}
 \item[\textbf{C1.}] supports et états ;
 \item[\textbf{C2.}] sélection interne ;
 \item[\textbf{C3.}] transporte;
 \item[\textbf{C4.}] lecteurs et quotients ;
 \item[\textbf{C5.}] mémoire et frontière ;
 \item[\textbf{C6.}] périodicité et retour ;
 \item[\textbf{C7.}] prolongement ;
 \item[\textbf{C8.}] salidas.
\end{enumerate}
La classe est déterminée par le domaine, le codomaine et la fonction, non par le nom
d’un chiffre associé.

\section{Quatorze regroupements fonctionnels}

\begingroup
\small
\renewcommand{\arraystretch}{1.12}
\begin{longtable}{@{}p{.20\textwidth}p{.29\textwidth}p{.43\textwidth}@{}}
\toprule
Regroupement&Opérateurs ou espaces&Fonction et invariants\\
\midrule
\endfirsthead
\toprule
Regroupement&Opérateurs ou espaces&Fonction et invariants\\
\midrule
\endhead
État et mémoire de route&
\(\mathcal X_{\rm TPK}^{\rm enr}\), \(\mathcal F_q\)&
Ils conservent la position, le feuillet, l’orientation, le résidu, le quotient, la retenue, la signature,
la frontière, le cylindre, le préfixe, la provenance et la mémoire.\\

Sélection interne&
\(M_{\rm ph}:\{1,\ldots,9\}\to\{-1,0,+1\}\)&
Active l’évaluation additive, multiplicative ou le pas de seuil ; ne déplace pas
l’état.\\

Transporte espacial&
\(T_N,T_E,T_S,T_O,F_{\rm loc},F_{\rm obs}\)&
Réalise des translations cardinales, une rétroaction de feuillet et une chronologie à deux
curseurs.\\

Lectores modulares&
\(\rho_9,q_9,\rho_3^{\rm or},c_3,q_{1000},r_{1000}\)&
Publient la phase, l’orientation et le bloc décimal en conservant les quotients de
relèvement.\\

Bloc et mémoire&
\(\mathcal K_{27}=F_{\rm obs}^{27}\), \(\mathcal N_{27}\)&
Le premier compose vingt-sept transports ; le second filtre des événements de
mémoire. Un indice identique n’implique pas un opérateur identique.\\

État dodécaphasique&
\(\mathcal R_{12},C_{12},\Theta_{108},\Gamma_{108}\)&
Distingue le registre enrichi, le secteur, la loi cyclique et le support ordonné.\\

Périodicité et profondeurs&
\(C_{12}\hookrightarrow C_{108}\twoheadrightarrow C_9\),
\(w_{108},w_{1080}\)&
Conserve l’extension non scindée, les mots de longueurs différentes et
la mémoire verticale.\\

Canaux \(90/120\)&
\(\mathcal F_6,\mathcal F_8,\mathcal I_{6\to8},
S_{90},S_{120}\)&
Sépare les incidences finies, les queues analytiques et leurs normalisations.\\

Bidirectionnalité&
\(P_\pm,\operatorname{rev},\operatorname{Ext}_\pm,
N_\pm,\beta,C_M\)&
Réalise l’échange, l’inversion de route, le prolongement futur et passé,
la valence, le bilan et la conjugaison.\\

Émission et coinduction&
\(L_0,L_1,R_{36},G_9,\varprojlim\operatorname{Cyl}_m\)&
Relève le germe, ouvre la frontière et organise des prolongements finis et
compatibles.\\

Lecteurs numériques&
\(\mathscr C_\pi,\mathscr P_e,F_{\rm av}\)&
Construisent des encadrements rationnels, une propagation factorielle et une auto-échelle de
Fibonacci.\\

Projection exceptionnelle&
\(\Gamma_{A_W}\), poids, support projectif&
Transporte la même émission vers une incidence ternaire sans sélectionner de chiffres
archimédiens.\\

Atlas et réalisation&
signature de route, atlas \(81/56/13/324\), \(C_M\)&
Classe les histoires et les caractères avant d’appliquer une carte dimensionnelle.\\

Validation finie&
recensements, ablations, empreintes de hachage et preuves d'indépendance&
Vérifie la couverture à la profondeur exécutée et distingue générateur, vérificateur et données de comparaison.\\
\bottomrule
\end{longtable}
\endgroup

Les huit classes décrivent la nature mathématique ; les quatorze regroupements,
la fonction d’exécution. Ce sont deux gradations du même inventaire.

\section{Registre nominal de trente-quatre entrées}

\begin{equation}
 \mathfrak O_{\rm TPK}:=(\mathfrak o_1,\ldots,\mathfrak o_{34}).
 \label{eq:u006f-registro-operadores}
\end{equation}

\begingroup
\scriptsize
\renewcommand{\arraystretch}{1.12}
\begin{longtable}{@{}r p{.19\textwidth}p{.28\textwidth}p{.42\textwidth}@{}}
\toprule
\#&Symbole&Domaine et codomaine&Règle ou fonction\\
\midrule
\endfirsthead
\toprule
\#&Symbole&Domaine et codomaine&Règle ou fonction\\
\midrule
\endhead
1&\(\mathcal A_9\)&\((\mathbb Z/9)^2\), matrice et loi interne&
Cellule arithmétique additive--multiplicative.\\
2&\((\rho_9,q_9)\)&\(\mathbb Z_{>0}\to
\{1,\ldots,9\}\times\mathbb Z_{\geq0}\)&
Division nonadique avec représentant positif et quotient.\\
3&\(\rho_3^{\rm or},c_3\)&\(\mathbb Z\to
\{-1,0,+1\}\times\mathbb Z\)&
Relèvement ternaire orienté \(n=\rho_3^{\rm or}(n)+3c_3(n)\).\\
4&\(\iota_{1000}:=(q_{1000},r_{1000})\)&\(\mathbb Z\to
\mathbb Z\times\{0,\ldots,999\}\)&
Division euclidienne \(N=1000q_{1000}+r_{1000}\).\\
5&\(T_d\)&\(X_9\to X_9\), \(d\in\{N,E,S,O\}\)&
Translation cardinale orientée.\\
6&\(F_{\rm loc}\)&\(\mathcal Q_{\rm loc}\times
\mathscr D_{\rm card}\to\mathcal Q_{\rm loc}\)&
Rétroaction locale de position et de feuillet.\\
7&\(T_{\min}\)&\(\Omega_{\rm seed}\to\mathcal U_6\)&
Émetteur minimal des deux curseurs, y compris la signature coordonnée.\\
8&\(G\)&\(\operatorname{im}s_\Sigma\times
\operatorname{im}s_\Pi\to\{0,\ldots,999\}^6\)&
Couplage décimal coordonnée par coordonnée.\\
9&\(\Psi_6\)&\(\mathcal U_6\to\mathbb F_3^6\)&
Réduction ternaire du catalogue.\\
10&\(\mathcal R_3\)&\(\mathcal Q_{\rm loc}\times
\mathscr D_{\rm card}^3\to\{-1,0,+1\}\)&
Facteur observable local de trois pas.\\
11&\(\mathcal K_{27}\)&\(\mathcal Q_{\rm obs}\to\mathcal Q_{\rm obs}\)&
Composition \(F_{\rm obs}^{27}\).\\
12&\(\mathcal X_{\rm TPK}^{\rm enr}\)&catégorie des états enrichis&
Conserve toutes les coordonnées typées du noyau.\\
13&\(\mathcal H\)&relation sur les états enrichis&
Transition de Hensel qui retient le résidu et le quotient.\\
14&\(\mathcal C_{3,1000}\)&états croisés
\((N,L,D,K,A,B)\)&
Actualisation exacte des cylindres ternaires--décimaux.\\
15&\(\Sigma_{\rm B}\)&mémoires de frontière \(\to\) signature conjointe&
Exprime l'orientation du bloc en conservant son registre.\\
16&\({\rm EF}\text{--}G_9\)&relation d’extension&
Filtre les extensions compatibles et organise la généalogie nonadique.\\
17&\(\Gamma_9\)&histoires enrichies \(\longrightarrow\) histoires&
Holonomie d’un tour de phase avec mémoire de frontière.\\
18&\(X_\chi\)&\(\varprojlim_m\operatorname{Surv}^{(\chi)}_m\)&
Section globale lorsque la non-vacuité, la compatibilité et l’unicité sont
démontrées à toute profondeur.\\
19&\(\mathcal P_{\rm APP}^{(2)}\)&catégorie bivariée de chemins&
Transporte simultanément l’évaluation multiplicative et l’accumulation
additive.\\
20&\(F_{\rm obs}\)&\(\mathcal Q_{\rm obs}\to\mathcal Q_{\rm obs}\)&
Chronologie observable déterministe à deux curseurs.\\
21&\(\mathcal A_{42}\)&alphabet de quarante-deux triades&
Support décimal élargi de l’émission finie.\\
22&\(\operatorname{Ext}_r\)&
\(\operatorname{Ext}_r\subseteq\mathcal F_q\times\mathcal F_{q'}\)&
Relation typée de prolongement des fibres.\\
23&\(\mathcal K_{\rm ph}\)&\(\mathcal U_6\times C_9\) sur lui-même&
Transfert de continuations compatibles avec phase explicite.\\
24&\((P,H,Q)\)&\(\mathcal X_\times\to\mathcal X_\times\)&
Avancée, demi-tour et quart de tour des germes.\\
25&\((c_{\rm mode},c_{\rm or})\)&caminos \(\to\mathbb Z^2\)&
Cocycles de changement de mode et d’orientation.\\
26&\(E_{20}\)&états finis \(\to\) vingt blocs&
Prolongement triadique fini avec livre de provenance.\\
27&\(\mathcal R_{12}\)&histoire enrichie \(\to\) registre de douze
secteurs&
Système temporel orienté avec parcours, signature, retenue et registre.\\
28&\(R_{\rm sgn}\)&état terminal \(\to U_{12}^{\rm sgn}
\subseteq\mathbb Z^{12}\)&
Lecteur de l’observable harmonique signée.\\
29&\(\mathcal H_{12}\)&\(\mathbb Z^{12}\to\mathbb Z^{12}\)&
Transformation par trois orbites de pas trois ; inverse
\(\mathcal H_{12}/4\) sur l’image pertinente.\\
30&\(\mathcal E_{\rm dod}=(D_3,D_4,Q)\)&
\(\mathbb Z^{12}\to\mathbb Z^{25}\)&
Extracteur transversal conjoint d’arités trois et quatre.\\
31&\(\mathcal C_\alpha\)&blocs et retenues \(\to\) blocs&
Récurrence dodécaphasique de clôture en base mille.\\
32&\((L_{\rm tick},\chi_{\rm mass})\)&histoire enrichie \(\to\)
caractère temporel&
Interface dimensionnelle ultérieure ; elle ne participe pas au constructeur numérique de
cet ouvrage.\\
33&\(\mathcal W_\pm\)&réseau bidirectionnel d’indice douze&
Lecture dans les deux sens en conservant l’orientation et la provenance.\\
34&\(\mathcal A_{\rm species}\)&familles d’extension \(\to\) atlas&
Interface de réalisation ultérieure, non un quatrième primitif.\\
\bottomrule
\end{longtable}
\endgroup

\section{Chaîne causale d'actualisation}

\begin{definition}[État canonique complet]
\label{def:u006f-esquema-estado}
Soit
\[
 \mathcal E:=
 \mathbb Z_{\rm quo}\times\mathbb Z_{\rm car}\times
 \mathsf{Mem}\times\mathsf{Term}\times\mathsf{Cyl}\times
 \mathsf{Fr}\times\mathsf{Pref}\times\Lambda .
\]
L’état enrichi canonique à l’instant \(t\) est
\begin{equation}
 x_t=(p_t,m_t,\sigma_t,b_t,r_t;
 q_t,c_t,\eta_t,s_t,C_t,\partial_t,N_t,\lambda_t)
 \in
 X_9\times\{\Sigma,\Pi\}\times\mathsf{TRIT}\times
 C_{12}\times C_9\times\mathcal E.
 \label{eq:u006f-esquema-estado}
\end{equation}
Les abréviations d’état utilisées dans d’autres unités sont des projections de
ce tuple. En particulier, ni le quotient \(q_t\), ni la retenue
\(c_t\), ni le terminal \(s_t\), ni le cylindre \(C_t\), ni la frontière
\(\partial_t\), ni le préfixe \(N_t\), ni le livre \(\lambda_t\) ne sont omis.
\end{definition}

Soit \(\rho_t\) la flèche observable déterminée par \(x_t\). Les lois de transport de U005 et la transition \(\mathcal C_{3,1000}\) déterminent un endomorphisme de mémoire. Pour \(z=(q,c,\eta,s,C,\partial,N,\lambda)\in\mathcal E\), on définit
\begin{align}
 \mathbf E_t:\mathcal E&\longrightarrow\mathcal E,\notag\\
 \mathbf E_t(z)&=\bigl(
 \mathsf H(\rho_t)q,
 \mathsf C(\rho_t)c+\kappa(\rho_t),
 \mathsf M(\rho_t)\eta,
 \mathsf T(\rho_t)s,
 \mathsf{Cyl}(\rho_t)C,
 \mathsf{Fr}(\rho_t)\partial,
 \mathsf{Pref}(\rho_t)N,
 \Lambda(\rho_t)\lambda
 \bigr).
 \label{eq:u006f-actualizacion-memoria}
\end{align}
Chaque composante publiée dans \eqref{eq:u006f-actualizacion-memoria} possède le
domaine indiqué par sa fibre ; \(\kappa(\rho_t)\) est la cochaîne de retenue.

Nous définissons trois endomorphismes du produit canonique. Si
\(\sigma'=M_{\rm ph}(g(t))\) et \(m'=\mu(\sigma',m)\), alors
\begin{align}
 \mathsf{Sel}_t(p,m,\sigma,b,r;e)
  &:=(p,m',\sigma',b,r;e),\label{eq:u006f-selector-tipado}\\
 \mathsf{Tra}_t(p,m,\sigma,b,r;e)
  &:=(T_{d(t)}p,m,\sigma,b,r;e),\label{eq:u006f-transporte-tipado}\\
 \mathsf{Upd}_t(p,m,\sigma,b,r;e)
  &:=(p,m,\sigma,b',r';\mathbf E_t(e)),
 \label{eq:u006f-memoria-tipada}
\end{align}
où \((b',r')\) est la cible de \(\rho_t:(b,r)\to(b',r')\). La chaîne causale est une composition d'applications, non la composition d'une application avec une valeur tritique :
\begin{equation}
 \boxed{\mathcal U_t:=
 \mathsf{Upd}_t\circ\mathsf{Tra}_t\circ\mathsf{Sel}_t.}
 \label{eq:u006f-tuberia}
\end{equation}

\begin{theorem}[Clôture de la chaîne]
\label{thm:u006f-cierre-tuberia}
L’application \(\mathcal U_t\) est un endomorphisme de l’état canonique
complet. Les lecteurs, prolongements et projections sont appliqués après
\(\mathcal U_t\) uniquement lorsque leur domaine a été produit.
\end{theorem}

\begin{proof}
\(\mathsf{Sel}_t\) ne change que le mode et le TRIT ;
\(\mathsf{Tra}_t\) ne change que la position APP ; et
\(\mathsf{Upd}_t\) applique le produit typé
\eqref{eq:u006f-actualizacion-memoria} et la flèche de phase. Les trois
codomaines sont le même produit de la
\cref{def:u006f-esquema-estado} ; leur composition est donc définie et
conserve toutes les coordonnées de l’état.
\end{proof}

\section{Règles de non-identification}

L’inventaire impose
\begin{equation}
 \mathcal S_4\neq D_4\neq\operatorname{sd}_4,
 \qquad
 \mathcal R_{12}\neq C_{12}\neq\Theta_{108}\neq\Gamma_{108},
 \qquad
 F_{\rm obs}\neq\mathcal K_{\rm ph},
 \label{eq:u006f-no-identificacion-uno}
\end{equation}
et
\begin{equation}
 U_{12}^{\rm cat}\neq U_{12}^{\rm sgn},
 \qquad
 \mathcal E_{\rm dod}\neq D_{\rm raw}^{90,120}
 \neq\mathcal K_{6\to8}.
 \label{eq:u006f-no-identificacion-dos}
\end{equation}
Ce sont des incompatibilités de type, non de simples inégalités de valeur.

\begin{criterion}[Audit de complétude opératorielle]
\label{crit:u006f-completitud}
Une exécution n’est opératoriellement complète que si :
\begin{enumerate}
 \item elle déclare les entrées de \(\mathfrak O_{\rm TPK}\) employées ;
 \item elle enregistre le domaine et le codomaine de chaque composition ;
 \item elle conserve l’ordre causal de \eqref{eq:u006f-tuberia} ;
 \item elle distingue dynamique, transfert, lecture, prolongement et sortie ;
 \item elle permet de reconstruire à partir de \(\lambda_t\) chaque changement de bloc et
 chaque retenue ;
 \item elle ne promeut pas une vérification finie en une section infinie sans
 prouver la non-vacuité, la compatibilité et l’unicité à toute profondeur.
\end{enumerate}
Tout symbole utilisé sans type ou toute fusion interdite par
\eqref{eq:u006f-no-identificacion-uno} et
\eqref{eq:u006f-no-identificacion-dos} constitue un échec vérifiable.
\end{criterion}
